% sec2_9.tex — Section 2.9: Least Squares Projection

\section{Least Squares Projection}
\label{sec:2.9}

What if the assumptions justifying the large-sample properties of the OLS estimator
(except for ergodic stationarity) are not satisfied but we nevertheless go ahead and
apply OLS to the sample? What is it that we estimate? The answer is in this section.
OLS provides an estimate of the best way to linearly combine the explanatory
variables to predict the dependent variable. The linear combination is called
the \textbf{least squares projection}.

\subsection*{Optimally Predicting the Value of the Dependent Variable}

We have been concerned about estimating unknown parameters from a sample.
Let us temporarily suspend the role of econometrician and put ourselves in the
following situation. There is a random scalar $y$ and a random vector $\mathbf{x}$. We know
the joint distribution of $(y, \mathbf{x})$ and the value of $\mathbf{x}$. On the basis of this knowledge
we wish to predict $y$. So a predictor is a function $f(\mathbf{x})$ of $\mathbf{x}$ with the functional
form $f(\cdot)$ determined by the joint distribution of $(y, \mathbf{x})$. Naturally, we choose the
function $f(\cdot)$ so as to minimize some index that is a function of the \textbf{forecast error}
$y - f(\mathbf{x})$. We take the loss function to be the \textbf{mean squared error} $\E[(y - f(\mathbf{x}))^2]$
because it seems as reasonable a loss function as any other and, more importantly,
because it produces the following convenient result:

\begin{proposition}[2.7]
\label{prop:2.7}
$\E(y \mid \mathbf{x})$ is the best predictor of $y$ in that it minimizes the mean
squared error.
\end{proposition}

We have seen in Chapter 1 that the add-and-subtract strategy is effective in showing
that the candidate solution minimizes a quadratic function. Let us apply the
strategy to the squared error here. Let $f(\mathbf{x})$ be any forecast. Add $\E(y \mid \mathbf{x})$ to the
forecast error $y - f(\mathbf{x})$ and then subtract it to obtain the decomposition
%
\begin{equation}
y - f(\mathbf{x}) = (y - \E(y \mid \mathbf{x})) + (\E(y \mid \mathbf{x}) - f(\mathbf{x})).
\label{eq:2.9.1}
\end{equation}
%
So the squared forecast error is
%
\begin{align}
(y - f(\mathbf{x}))^2 &= (y - \E(y \mid \mathbf{x}))^2
  + 2(y - \E(y \mid \mathbf{x}))(\E(y \mid \mathbf{x}) - f(\mathbf{x})) \notag \\
  &\quad + (\E(y \mid \mathbf{x}) - f(\mathbf{x}))^2.
\label{eq:2.9.2}
\end{align}
%
Take the expectation of both sides to obtain
%
\begin{align}
\text{mean squared error} &\equiv \E\!\big[(y - f(\mathbf{x}))^2\big] \notag \\
&= \E\!\big[(y - \E(y \mid \mathbf{x}))^2\big]
  + 2\,\E\!\big[(y - \E(y \mid \mathbf{x}))(\E(y \mid \mathbf{x}) - f(\mathbf{x}))\big] \notag \\
&\quad + \E\!\big[(\E(y \mid \mathbf{x}) - f(\mathbf{x}))^2\big].
\label{eq:2.9.3}
\end{align}
%
It is a straightforward application of the Law of Total Expectations to show that
the middle term, which is the covariance between the optimal forecast error and the
difference in forecasts, is zero (a review question). Therefore,
%
\begin{align}
\text{mean squared error} &= \E\!\big[(y - \E(y \mid \mathbf{x}))^2\big]
  + \E\!\big[(\E(y \mid \mathbf{x}) - f(\mathbf{x}))^2\big] \notag \\
&\geq \E\!\big[(y - \E(y \mid \mathbf{x}))^2\big],
\label{eq:2.9.4}
\end{align}
%
which shows that the mean squared error is bounded from below by
$\E[(y - \E(y \mid \mathbf{x}))^2]$, and this lower bound is achieved by the conditional expectation.

\subsection*{Best Linear Predictor}

It requires the knowledge of the joint distribution of $(y, \mathbf{x})$ to calculate
$\E(y \mid \mathbf{x})$, which may be highly nonlinear. We now restrict the predictor to being
a linear function of $\mathbf{x}$ and ask: what is the best (in the sense of minimizing the mean
squared error) linear predictor of $y$ based on $\mathbf{x}$? For this purpose, consider
$\bm{\beta}^*$ that satisfies the orthogonality condition
%
\begin{equation}
\E[\mathbf{x} \cdot (y - \mathbf{x}'\bm{\beta}^*)] = \mathbf{0}
\quad \text{or} \quad
\E(\mathbf{x}\mathbf{x}')\bm{\beta}^* = \E(\mathbf{x} \cdot y).
\label{eq:2.9.5}
\end{equation}
%
The idea is to choose $\bm{\beta}^*$ so that the forecast error $y - \mathbf{x}'\bm{\beta}^*$ is orthogonal to
$\mathbf{x}$. If $\E(\mathbf{x}\mathbf{x}')$ is nonsingular, the orthogonality condition can be solved for
$\bm{\beta}^*$:
%
\begin{equation}
\bm{\beta}^* = [\E(\mathbf{x}\mathbf{x}')]^{-1}\,\E(\mathbf{x} \cdot y).
\label{eq:2.9.6}
\end{equation}
%
The \textbf{least squares} (or \textbf{linear}) \textbf{projection} of $y$ on $\mathbf{x}$, denoted
$\widehat{\E}^*(y \mid \mathbf{x})$, is defined as $\mathbf{x}'\bm{\beta}^*$, where $\bm{\beta}^*$ satisfies
\eqref{eq:2.9.5} and is called the \textbf{least squares projection coefficients}.

\begin{proposition}[2.8]
\label{prop:2.8}
The least squares projection $\widehat{\E}^*(y \mid \mathbf{x})$ is the best linear predictor
of $y$ in that it minimizes the mean squared error.
\end{proposition}

The add-and-subtract strategy also works here.

\begin{proof}
For any linear predictor $\mathbf{x}'\widetilde{\bm{\beta}}$,
%
\begin{align*}
\text{mean squared error} &\equiv \E\!\big[(y - \mathbf{x}'\widetilde{\bm{\beta}})^2\big] \\
&= \E\!\big\{\big[(y - \mathbf{x}'\bm{\beta}^*)
  + \mathbf{x}'(\bm{\beta}^* - \widetilde{\bm{\beta}})\big]^2\big\}
  \quad \text{(by the add-and-subtract strategy)} \\
&= \E\!\big[(y - \mathbf{x}'\bm{\beta}^*)^2\big]
  + 2(\bm{\beta}^* - \widetilde{\bm{\beta}})'\,
  \E\!\big[\mathbf{x} \cdot (y - \mathbf{x}'\bm{\beta}^*)\big]
  + \E\!\big[(\mathbf{x}'(\bm{\beta}^* - \widetilde{\bm{\beta}}))^2\big] \\
&= \E\!\big[(y - \mathbf{x}'\bm{\beta}^*)^2\big]
  + \E\!\big[(\mathbf{x}'(\widetilde{\bm{\beta}} - \bm{\beta}^*))^2\big] \\
&\qquad \text{(by the orthogonality condition \eqref{eq:2.9.5})} \\
&\geq \E\!\big[(y - \mathbf{x}'\bm{\beta}^*)^2\big]. \qedhere
\end{align*}
\end{proof}

In contrast to the best predictor, which is the conditional expectation, the best linear
predictor requires only the knowledge of the second moments of the joint distribution of $(y, \mathbf{x})$ to calculate (see \eqref{eq:2.9.6}).

If one of the regressors $\mathbf{x}$ is a constant, the least squares projection coefficients
can be written in terms of variances and covariances. Let $\widetilde{\mathbf{x}}$ be the vector of non-constant regressors so that
%
\[
\mathbf{x} = \begin{bmatrix} 1 \\ \widetilde{\mathbf{x}} \end{bmatrix}.
\]
%
An analytical exercise to this chapter asks you to prove that
%
\begin{equation}
\widehat{\E}^*(y \mid \mathbf{x})
\equiv \widehat{\E}^*(y \mid 1, \widetilde{\mathbf{x}})
= \mu + \bm{\gamma}'\widetilde{\mathbf{x}},
\label{eq:2.9.7}
\end{equation}
%
where
\[
\bm{\gamma} = \Var(\widetilde{\mathbf{x}})^{-1}\,\Cov(\widetilde{\mathbf{x}}, y),
\quad
\mu = \E(y) - \bm{\gamma}'\,\E(\widetilde{\mathbf{x}}).
\]
%
This formula is the population analogue of the ``deviations-from-the-mean regression'' formula of Chapter~1 (Analytical Exercise~3).

\subsection*{OLS Consistently Estimates the Projection Coefficients}

Now let us put the econometrician's hat back on and consider estimating $\bm{\beta}^*$.
Suppose we have a sample of size $n$ drawn from an ergodic stationary stochastic
process $\{y_i, \mathbf{x}_i\}$ with the joint distribution $(y_i, \mathbf{x}_i)$ (which does not depend on $i$
because of stationarity) identical to that of $(y, \mathbf{x})$ above. So, for example,
$\E(\mathbf{x}_i\mathbf{x}_i') = \E(\mathbf{x}\mathbf{x}')$.
By the Ergodic Theorem the second moments in \eqref{eq:2.9.6} can be consistently estimated by the corresponding sample second moments. Thus a consistent estimator
of the projection coefficients $\bm{\beta}^*$ is
%
\[
\Bigl(\frac{1}{n}\sum_{i=1}^{n}\mathbf{x}_i\mathbf{x}_i'\Bigr)^{-1}
\Bigl(\frac{1}{n}\sum_{i=1}^{n}\mathbf{x}_i \cdot y_i\Bigr)
= (\mathbf{X}'\mathbf{X})^{-1}\mathbf{X}'\mathbf{y},
\]
%
which is none other than the OLS estimator $\mathbf{b}$. That is, under Assumption~2.2
(ergodic stationarity) and Assumption~2.4 guaranteeing the nonsingularity of
$\E(\mathbf{x}\mathbf{x}')$, the OLS estimator is always consistent for the projection coefficient vector,
the $\bm{\beta}^*$ that satisfies the orthogonality condition \eqref{eq:2.9.5}.

\bigskip
\noindent\textsc{Questions for Review}
\begin{enumerate}
\item (Unforecastability of forecast error)\quad
For the minimum mean square error forecast $\E(y \mid \mathbf{x})$, the forecast error is orthogonal to any function $\phi(\mathbf{x})$ of $\mathbf{x}$. That is, $\E[\eta\phi(\mathbf{x})] = 0$ where
$\eta \equiv y - \E(y \mid \mathbf{x})$. Prove this.
\textbf{Hint:} The Law of Total Expectations. Show that the middle term on the RHS of \eqref{eq:2.9.3} is zero by setting $\phi(\mathbf{x}) = \E(y \mid \mathbf{x}) - f(\mathbf{x})$.

\item (Forecasting white noise)\quad
Suppose $\{\varepsilon_i\}$ is white noise. What is
$\widehat{\E}^*(\varepsilon_i \mid \varepsilon_{i-1}, \varepsilon_{i-2}, \ldots, \varepsilon_{i-m})$?
What is
$\widehat{\E}^*(\varepsilon_i \mid 1, \varepsilon_{i-1}, \varepsilon_{i-2}, \ldots, \varepsilon_{i-m})$?
Is it true that $\E(\varepsilon_i \mid \varepsilon_{i-1}, \varepsilon_{i-2}, \ldots, \varepsilon_{i-m}) = 0$?
\textbf{Hint:} Is it zero for the process of Example~2.4?

\item (Conditional expectations that are linear)\quad
Suppose $\E(y \mid \widetilde{\mathbf{x}}) = \mu + \bm{\gamma}'\widetilde{\mathbf{x}}$. Show:
$\widehat{\E}^*(y \mid 1, \widetilde{\mathbf{x}}) = \E(y \mid \widetilde{\mathbf{x}})$.

\item (Partitioned projection)\quad
Consider the model $y_i = \mathbf{x}_i'\bm{\beta} + \mathbf{z}_i'\bm{\delta} + \varepsilon_i$
with $\E(\mathbf{x}_i \cdot \varepsilon_i) = \mathbf{0}$, $\E(\mathbf{z}_i \cdot \varepsilon_i) \neq \mathbf{0}$, and
$\E(\mathbf{z}_i\mathbf{x}_i') = \mathbf{0}$. Thus, $\mathbf{z}_i$ is not predetermined (i.e., not orthogonal to the error term), but it is unrelated to the predetermined regressor $\mathbf{x}_i$ in that the cross moments are zero.
\begin{enumerate}[label=(\alph*)]
\item Show that the least squares projection coefficient of $\mathbf{x}_i$ in the projection of
$y_i$ on $\mathbf{x}_i$ and $\mathbf{z}_i$ is $\bm{\beta}$.
\textbf{Hint:} Calculate $\widehat{\E}^*(\varepsilon_i \mid \mathbf{x}_i, \mathbf{z}_i)$.
\item What is the least squares projection coefficient of $\mathbf{x}_i$ in the least squares
projection of $y_i$ on $\mathbf{x}_i$?
\textbf{Hint:} Treat $\mathbf{z}_i'\bm{\delta} + \varepsilon_i$ as the error term.
\item Which projection would you use for estimating $\bm{\beta}$?
\textbf{Hint:} You want the error variance to be smaller.
\end{enumerate}
\end{enumerate}
